% Official tau2 environment validation (2026-08-25)
\section{Positive Control: Full Task Completion on the Official tau2 Benchmark}

The audit harness is deployed on the official \texttt{tau2} benchmark
environment to check that the serving runtime itself is competitive on a
real, externally-scored task. The setup uses the benchmark's own agent
(\texttt{llm\_agent}), user simulator (\texttt{user\_simulator}),
orchestrator, environment, and hidden evaluator --- including the LLM judge
for natural-language assertions --- untouched. The only substitution is the
model backend: instead of a commercial API, all LLM calls (agent, user, and
judge) are served by our policy-consistent local runtime
(\texttt{ColocatedPolicy}, Qwen2.5-14B-Instruct) through a small
OpenAI-compatible endpoint. This preserves the property the pipeline was
built for: the model that generates is the model that would receive
gradient updates, and no request ever leaves the machine.

Serving-side engineering (no environment or evaluator modification) was
required to make a 14B model competitive on the 16-tool retail domain:

\begin{itemize}
\item \textbf{Tool schema injection}: the official agent's system prompt is
  augmented with the actual tool list (names, argument types, one-line
  descriptions) and a short usage policy (never invent ids; count options
  by listing product types, then counting \emph{available} variants; look
  up order details before any mutation; confirm with the user before
  changes).
\item \textbf{Tool-result digests}: large JSON tool outputs (order details,
  user details, product variants) are rendered to the model as compact
  line-oriented summaries. The model reads \texttt{item 4602305039: Vacuum
  Cleaner (\$561.05)} instead of a nested JSON blob; a 14B model otherwise
  misreads items inside multi-hundred-byte payloads.
\item \textbf{Anti-loop valve}: a tool call identical to the last executed
  call is never returned; the agent is forced to change strategy (or ask
  the user) instead of repeating a failing lookup.
\end{itemize}

On the official retail task set, the agent completes both evaluated tasks
end-to-end with full reward in 5/10 seeded runs (2/5 on the exchange
task, 3/5 on the count-and-return task):

\begin{center}
\begin{tabular}{lccccc}
\toprule
task & seed 0 & seed 1 & seed 2 & seed 3 & seed 4 \\
\midrule
0 (exchange 2 items)       & 1.0 & 0.0 & 0.0 & 1.0 & 0.0 \\
2 (count + return 3 items) & 1.0 & 1.0 & 0.0 & 0.0 & 1.0 \\
\bottomrule
\end{tabular}
\end{center}

Each reward is the product of the task's reward-basis components (DB
state change and the LLM-judged NL assertion), so $1.0$ means the
environment mutation was exactly as required \emph{and} the agent
communicated the required information. Failures are last-mile
single-call errors --- one item omitted from a multi-item return call,
or a fabricated payment-method id late in an otherwise-correct exchange
--- a model-capability limit rather than a pipeline defect: the
successful trajectories and the failing ones execute the same full flow
(user lookup, product counting, order location).

The task-2 trajectory is a full solution: find the user by name and zip;
list product types; fetch T-shirt details and report ``there are 10
t-shirt options available'' (the hidden judge independently confirms the
statement); walk the user's five orders to locate the items; and execute
\texttt{return\_delivered\_order\_items} with the three real item ids and
the real payment method. The environment-side evaluation confirms the
database was mutated exactly as required.

This is the positive control the audit arc needed: with a clean serving
runtime and modest prompt/serving engineering, the same harness that
detects silent failure modes in agent test-time RL also \emph{solves} a
real benchmark task end-to-end with local models and no external API. It
separates ``the pipeline cannot do the task'' (false --- task 0 and 2
reward 1.0) from ``the pipeline cannot safely learn'' (true under the
clean protocol, F1--F3): the failure to transfer is attributable to the
update rule and the validation distribution, not to the serving or
evaluation machinery.

\subsection{The audit arc replicates on the official environment}
The same stream machinery that produced the F3 finding on the in-house
CTS stream was run on the official tau2 environment: a 10-task
prequential stream (retail base tasks 0--9) through the official
orchestrator, with accessible-evidence credit only (a task is a success
iff the trajectory contains a mutation call whose id-valued arguments
all appear in the tool results of the same conversation) and the hidden
official evaluator's reward recorded for reporting only. Two arms,
common random numbers across arms:

\begin{center}
\begin{tabular}{lcc}
\toprule
arm & accessible success (seed 0 / seed 1) & commits \\
\midrule
frozen & 1/10 (task 2) / 1/10 (task 5) & --- \\
pair   & 1/10 (task 2) & 0/6 (all gate-rejected) \\
\bottomrule
\end{tabular}
\end{center}

The pair arm trained a shadow candidate after every task that yielded
positive and negative rows (six updates over the stream) and submitted
it to the pre-commit gate; every candidate scored a $0.00$ improvement
rate on the validation prompt and was discarded. The served policy never
changed, so the pair arm's trajectory is bit-identical to frozen's on
seed 0 (including the same task-2 success). The gate's rejection rate is
the tau2 replication of the v3.2 finding: on this environment the
candidate's gain does not generalize from the validation instances to
the production distribution, so the protective gate is a safe no-op and
agent test-time RL still shows no transfer under the clean protocol.
Combined with the F3 result on the in-house stream, the audit arc's
conclusion now holds across two environments, one of them externally
scored with an official hidden evaluator.
